% Generated draft: requires mathematical and source review.
\documentclass[11pt]{article}
\usepackage[margin=25mm]{geometry}
\usepackage{fontspec}
\setmainfont{Latin Modern Roman}
\usepackage{amsmath,amssymb,mathtools,unicode-math}
\setmathfont{Latin Modern Math}
\usepackage{enumitem,needspace,xurl,hyperref}
\setlength{\parindent}{0pt}
\setlength{\parskip}{5pt}
\setlength{\emergencystretch}{3em}
\newcommand{\N}{\mathbb{N}}
\newcommand{\Z}{\mathbb{Z}}
\newcommand{\Q}{\mathbb{Q}}
\newcommand{\R}{\mathbb{R}}
\newcommand{\C}{\mathbb{C}}
\newcommand{\F}{\mathbb{F}}
\newcommand{\A}{\mathbb{A}}
\newcommand{\PP}{\mathbb{P}}
\newcommand{\E}{\mathbb{E}}
\newcommand{\eps}{\varepsilon}
\newcommand{\dd}{\,\mathrm{d}}
\begin{document}
\section*{UnsolvedMath ID: 30004646}
\textbf{Source problem code:} catalogue entry
\subsection{Definitions and mathematical statement}
Let \(X\) be a connected smooth quasi-projective variety over \(\C\), let \(x\in X(\C)\), and fix a smooth projective compactification \(X\hookrightarrow\overline X\) whose boundary \(D=\overline X\setminus X\) is a normal-crossings divisor. Let \(n\geq 1\) be an integer and let
\[
\rho:\pi_1^{\mathrm{top}}(X(\C),x)\longrightarrow\operatorname{GL}_n(\C)
\]
be a representation. Write \(\mathcal L_\rho\) for its associated rank-\(n\) complex local system on \(X(\C)\). Irreducibility means that \(\C^n\) has no nonzero proper \(\rho\)-invariant subspace. Quasi-unipotence along the boundary means that, for every irreducible component of \(D\), the monodromy \(T\) around a small loop about a general smooth point of that component satisfies \(T^k\) unipotent for some integer \(k\geq 1\), which may depend on the component.

Assume also that \(\rho\) is rigid. The supplied record does not define rigidity or specify the deformation problem in which rigidity is imposed. In particular, it does not say whether boundary monodromy conjugacy classes are fixed, whether rigidity is local isolation or global determination by boundary data, or whether cohomological rigidity is intended. Consequently, the following question retains this unresolved hypothesis and requires a choice of its intended meaning before it becomes a fully precise assertion.

\emph{Question under the source's rigidity hypothesis.} For every \(X\), compactification, and \(\rho\) satisfying the preceding assumptions, do there exist a nonempty dense Zariski-open subset \(U\subseteq X\), an integer \(m\geq 0\), a smooth projective morphism of complex algebraic varieties \(f:Y\to U\), and a complex local system \(\mathcal M\) on \(U(\C)\) such that
\[
R^m(f^{\mathrm{an}})_*\underline{\C}_{Y(\C)}
\cong
\mathcal L_\rho|_{U(\C)}\oplus\mathcal M
\]
as complex local systems on \(U(\C)\)? Here \(f^{\mathrm{an}}:Y(\C)\to U(\C)\) is the associated analytic map, and \(\underline{\C}_{Y(\C)}\) is the constant sheaf. Equivalently, after choosing a base point in \(U(\C)\), the restricted representation is required to be a direct summand of the monodromy representation on the degree-\(m\) cohomology of the fibers of \(f\). The open subset, degree, and family may depend on \(\rho\). No determinant condition is supplied in the record, so none is imposed here.
\subsection{Short English statement}
The question asks whether every irreducible rigid complex local system with quasi-unipotent boundary monodromy on a connected smooth quasi-projective complex variety becomes a direct summand of the cohomology local system of a smooth projective family after restriction to a dense open subset. The intended meaning of rigidity remains unspecified. The supplied partially solved label and literature assessment are dated curation, not independently verified conclusions.
\subsection{Sources}
\begin{itemize}
\item Local UnsolvedMath problem registry, permanent ID \(30004646\), exact source problem code \texttt{catalogue entry}, title \emph{Motivic Realization of Rigid Local Systems}, in the Oberwolfach Reports Open Problems collection. The supplied record gives proposed year 2021 and identifies no individual proposer.
\item Supplied report citation: \emph{Homotopic and Geometric Galois Theory (2022)}, Oberwolfach Reports, pp.\ 663--744, DOI \texttt{10.4171/owr/2021/12}. \url{https://doi.org/10.4171/owr/2021/12}. The literature-source list labels this same reference \emph{Oberwolfach Report 12/2021}. These supplied year designations are retained without reconciliation.
\item Supplied literature source: \emph{Motivic realization of rigid G-local systems on curves}. \url{https://arxiv.org/abs/2405.18268}. The supplied assessment attributes a motivic realization result for rigid reductive-group local systems on curves to F\ae rgeman; the precise hypotheses and relationship to the formulation above have not been independently verified.
\item Supplied literature source: \emph{Cohomologically rigid local systems and integrality}. \url{https://arxiv.org/abs/1711.06436}.
\item Supplied literature assessment, checked on exactly \texttt{2026-08-22}: the record labels the problem partially solved, states that Katz and later work settle important curve and cohomologically rigid cases, attributes the stated curve realization result to F\ae rgeman, and states that the general smooth quasi-projective variety case remains open. This is supplied dated curation derived from the OWR record, not an independently verified literature assessment. The supplied background expressly requests independent verification of time-sensitive status and cited claims.
\end{itemize}
\end{document}
